\begin{proof}[Proof of \cref{obj:thm:finite-oracle}]
Fix the parameters satisfying the stated assumptions. Throughout the proof, write
\[
c=\begin{pmatrix}-1\\1\end{pmatrix},
\qquad
v(t)=\begin{pmatrix}t\\t+1\end{pmatrix},
\qquad
F_\theta(\alpha,t)
=
\sum_{S\in\mathcal S}
\alpha_S\frac{(g_S^\top v(t))^2}{h_S(\theta)}
\quad
(\alpha\in\mathbb R^{\mathcal S},\ t\in\mathbb R).
\]
Thus \(c^\top v(t)=1\). The staircase quantities are those in
\cref{obj:lem:staircase-refinement}, and the optimized values are those in
\cref{obj:def:staircase-saddle}.

\begin{enumerate}
\item We first establish the existence facts needed for the minimax argument.
For \(S\in\mathcal S\) and \(x\in\mathcal X\), define
\[
b_S(x)=
\begin{cases}
e^\varepsilon,&x\in S,\\
1,&x\notin S.
\end{cases}
\]
Throughout this proof, relabel the four records indexed \(1,2,3,4\)
in \cref{obj:lem:staircase-refinement} by \(0,1,2,3\), respectively,
preserving their order and reindexing their probability coordinates
accordingly. The record probabilities are therefore
\[
(\pi_{\theta,0},\pi_{\theta,1},\pi_{\theta,2},\pi_{\theta,3})
=
\bigl(
(1-p)(1-\mu_0),\ (1-p)\mu_0,\
p(1-\mu_1),\ p\mu_1
\bigr).
\]
They are positive, and hence
\[
h_S(\theta)=\sum_{j=0}^3\pi_{\theta,j}b_S(x_j)>0
\qquad(S\in\mathcal S).
\]
Let the complementary-pattern weights be
\[
(\alpha_{\mathrm{pair}})_S=
\begin{cases}
(e^\varepsilon+1)^{-1},
&S=\{x_0\}\ \text{or}\ S=\mathcal X\setminus\{x_0\},\\
0,&\text{otherwise}.
\end{cases}
\]
Each record belongs to exactly one of these two patterns, so
\[
\sum_{S\in\mathcal S}
(\alpha_{\mathrm{pair}})_S b_S(x)
=\frac{e^\varepsilon+1}{e^\varepsilon+1}=1
\qquad(x\in\mathcal X).
\]
Thus \(\mathcal A_\varepsilon\) is nonempty. Its defining inequalities
and equalities are closed and convex. Moreover, since \(b_S(x)\ge1\),
\[
0\le\alpha_S\le\sum_{T\in\mathcal S}\alpha_T
\le\sum_{T\in\mathcal S}\alpha_Tb_T(x)=1
\qquad
(\alpha\in\mathcal A_\varepsilon,\ S\in\mathcal S).
\]
Consequently, \(\mathcal A_\varepsilon\) is compact.

For every feasible \(\alpha\), define the quadratic coefficients
\[
a_{\mathrm{quad}}(\alpha)
=
2\sum_{S\in\mathcal S}
\alpha_S
\frac{(g_S^\top v(0))(g_{S,0}+g_{S,1})}{h_S(\theta)},
\qquad
b_{\mathrm{quad}}(\alpha)
=
\sum_{S\in\mathcal S}
\alpha_S
\frac{(g_{S,0}+g_{S,1})^2}{h_S(\theta)}\ge0,
\]
where \(g_S=(g_{S,0},g_{S,1})^\top\). Expanding the squares gives
\[
F_\theta(\alpha,u)
=
F_\theta(\alpha,0)
+u\,a_{\mathrm{quad}}(\alpha)
+u^2b_{\mathrm{quad}}(\alpha)
\qquad(u\in\mathbb R).
\]
If \(b_{\mathrm{quad}}(\alpha)=0\), nonnegativity of
\(F_\theta(\alpha,u)\) for every \(u\) forces
\(a_{\mathrm{quad}}(\alpha)=0\): otherwise,
\[
F_\theta\!\left(
\alpha,-\frac{F_\theta(\alpha,0)+1}{a_{\mathrm{quad}}(\alpha)}
\right)=-1.
\]
If \(b_{\mathrm{quad}}(\alpha)>0\), completing the square gives a
minimizer. In both cases, a minimizing coordinate is
\[
t_\alpha=
\begin{cases}
0,&b_{\mathrm{quad}}(\alpha)=0,\\[2mm]
-\dfrac{a_{\mathrm{quad}}(\alpha)}
{2b_{\mathrm{quad}}(\alpha)},
&b_{\mathrm{quad}}(\alpha)>0,
\end{cases}
\]
and
\[
F_\theta(\alpha,u)-F_\theta(\alpha,t_\alpha)
=
b_{\mathrm{quad}}(\alpha)(u-t_\alpha)^2\ge0.
\]
In particular, the profile
\[
\mathscr P(\alpha)=\min_{t\in\mathbb R}F_\theta(\alpha,t)
\qquad(\alpha\in\mathcal A_\varepsilon)
\]
is finite and nonnegative. For each fixed \(t\), the objective is
continuous in \(\alpha\); its infimum over \(t\) is therefore upper
semicontinuous. Compactness supplies a feasible
\(\alpha_{\max}\) such that
\[
\mathscr P(\alpha)\le\mathscr P(\alpha_{\max})
\qquad(\alpha\in\mathcal A_\varepsilon),
\qquad
J^*(\theta,p,\varepsilon)=\mathscr P(\alpha_{\max}).
\]

Define the upper envelope by
\[
\mathscr E(t)=
\max_{\alpha\in\mathcal A_\varepsilon}F_\theta(\alpha,t)
\qquad(t\in\mathbb R).
\]
The maximum exists by compactness and continuity. Joint continuity
of \(F_\theta\), together with compactness of the weight set, also
makes \(\mathscr E\) continuous.

For the first singleton pattern,
\[
g_{\{x_0\}}
=
\begin{pmatrix}-(e^\varepsilon-1)(1-p)\\0\end{pmatrix}.
\]
Define the positive envelope coefficient
\[
\eta_{\mathrm{env}}
=
\frac{(e^\varepsilon-1)^2(1-p)^2}
{(e^\varepsilon+1)h_{\{x_0\}}(\theta)}>0.
\]
Nonnegativity of every summand yields
\[
\eta_{\mathrm{env}}t^2
\le F_\theta(\alpha_{\mathrm{pair}},t)
\le\mathscr E(t)
\qquad(t\in\mathbb R).
\]
Hence the sublevel set at \(\mathscr E(0)\) satisfies
\[
\{t\in\mathbb R:\mathscr E(t)\le\mathscr E(0)\}
\subseteq
\left[
-\sqrt{\frac{\mathscr E(0)}{\eta_{\mathrm{env}}}},
\sqrt{\frac{\mathscr E(0)}{\eta_{\mathrm{env}}}}
\right].
\]
It is nonempty and closed, and therefore compact. Minimizing on
this set gives a coordinate \(t_{\mathrm{env}}\in\mathbb R\) with
\[
\mathscr E(t_{\mathrm{env}})
\le\mathscr E(t)
\qquad(t\in\mathbb R).
\]
% lean: CausalSmith.Stat.LdpAteEfficiencySurface.informationObjective_exists_minimizer
% lean: CausalSmith.Stat.LdpAteEfficiencySurface.informationObjective_profile_exists_maximizer
% lean: CausalSmith.Stat.LdpAteEfficiencySurface.upperEnvelope_exists_minimizer

\item We now prove the minimax identity. Define
\[
\mathscr N(\alpha,t)=-F_\theta(\alpha,t)
\qquad
(\alpha\in\mathcal A_\varepsilon,\ t\in\mathbb R).
\]
For fixed \(t\), this function is continuous and affine in
\(\alpha\), hence quasiconvex. For fixed feasible \(\alpha\), it is
continuous and concave in \(t\), since
\(b_{\mathrm{quad}}(\alpha)\ge0\), hence quasiconcave. The weight
set is nonempty, compact, and convex, and \(\mathbb R\) is convex.
The extrema established above identify the quantities entering
Sion's minimax theorem \citep{Sion1958}:
\[
\sup_{t\in\mathbb R}\mathscr N(\alpha,t)=-\mathscr P(\alpha),
\qquad
\inf_{\alpha\in\mathcal A_\varepsilon}
\sup_{t\in\mathbb R}\mathscr N(\alpha,t)
=-J^*(\theta,p,\varepsilon),
\]
and
\[
\inf_{\alpha\in\mathcal A_\varepsilon}\mathscr N(\alpha,t)
=-\mathscr E(t),
\qquad
\sup_{t\in\mathbb R}
\inf_{\alpha\in\mathcal A_\varepsilon}\mathscr N(\alpha,t)
=-\inf_{t\in\mathbb R}\mathscr E(t).
\]
Sion's theorem equates the last quantities in these two displays.
Changing signs gives
\[
J^*(\theta,p,\varepsilon)
=
\inf_{t\in\mathbb R}\mathscr E(t)
=
\inf_{t\in\mathbb R}
\max_{\alpha\in\mathcal A_\varepsilon}F_\theta(\alpha,t).
\]
% lean: CausalSmith.Stat.LdpAteEfficiencySurface.finiteOracle_minimax

\item For fixed \(t\), define the vertex envelope
\[
\mathscr E_{\mathrm{vert}}(t)
=
\sup_{\alpha\in\operatorname{vert}(\mathcal A_\varepsilon)}
F_\theta(\alpha,t).
\]
The compact convex set \(\mathcal A_\varepsilon\) has extreme
points, and its continuous objective is bounded above there.
Since its vertices belong to the feasible set,
\[
\mathscr E_{\mathrm{vert}}(t)\le\mathscr E(t).
\]
Conversely, define the sublevel set
\[
\mathcal C_t=
\{\alpha\in\mathbb R^{\mathcal S}:
F_\theta(\alpha,t)\le\mathscr E_{\mathrm{vert}}(t)\}.
\]
It is closed by continuity and convex by affinity of the objective.
It contains all extreme points. The closed-convex-hull
characterization of a compact convex set therefore gives
\[
\mathcal A_\varepsilon
=
\overline{\operatorname{conv}
\bigl(\operatorname{vert}(\mathcal A_\varepsilon)\bigr)}
\subseteq\mathcal C_t.
\]
Evaluating at a maximizing feasible weight proves
\[
\mathscr E(t)\le\mathscr E_{\mathrm{vert}}(t),
\qquad
\mathscr E(t)=\mathscr E_{\mathrm{vert}}(t).
\]
The feasible set is a bounded polyhedron in
\(\mathbb R^{\mathcal S}\), so its vertex set is finite and the
vertex supremum is a maximum. Thus the two envelope value sets
coincide, and
\[
J^*(\theta,p,\varepsilon)
=
\inf_{t\in\mathbb R}
\max_{\alpha\in\operatorname{vert}(\mathcal A_\varepsilon)}
F_\theta(\alpha,t).
\]
% lean: CausalSmith.Stat.LdpAteEfficiencySurface.upperEnvelope_eq_vertexEnvelope

\item Let \(Q\) be any admissible stationary channel on an arbitrary
measurable output space. By \cref{obj:lem:staircase-refinement},
there are feasible weights \(\alpha_Q\) and a probability kernel
\(K_Q\) such that
\[
Q=K_Q\circ Q_{\alpha_Q},
\qquad
I_\theta(\alpha_Q)-I_\theta(Q)\succeq0.
\]
Choose a minimizing coordinate \(t_Q\in\mathbb R\) for the
quadratic objective of \(\alpha_Q\). The profile maximum established
above gives
\[
F_\theta(\alpha_Q,t_Q)
=\mathscr P(\alpha_Q)
\le\mathscr P(\alpha_{\max})
=J^*(\theta,p,\varepsilon).
\]
Expanding the information matrix shows, for every feasible
\(\alpha\) and every \(t\in\mathbb R\), that
\[
v(t)^\top I_\theta(\alpha)v(t)
=
\sum_{S\in\mathcal S}
\alpha_S\frac{(g_S^\top v(t))^2}{h_S(\theta)}
=F_\theta(\alpha,t).
\]
Applying the positive semidefinite difference to \(v(t_Q)\) gives
the intermediate bound
\[
v(t_Q)^\top I_\theta(Q)v(t_Q)
\le
v(t_Q)^\top I_\theta(\alpha_Q)v(t_Q)
=
F_\theta(\alpha_Q,t_Q)
\le J^*(\theta,p,\varepsilon).
\]

We next verify that the quantity on the right is positive. Define
the singleton weights
\[
(\alpha_{\mathrm{sing}})_S=
\begin{cases}
(e^\varepsilon+3)^{-1},&S=\{x_j\}\ \text{for some }j\in\{0,1,2,3\},\\
0,&\text{otherwise}.
\end{cases}
\]
For each record, one singleton contributes \(e^\varepsilon\)
and the other three contribute \(1\), so
\[
\sum_{S\in\mathcal S}
(\alpha_{\mathrm{sing}})_S b_S(x)
=\frac{e^\varepsilon+3}{e^\varepsilon+3}=1.
\]
These weights are feasible. Define the positive coefficients
\[
\eta_{\mathrm{control}}
=
\frac{(e^\varepsilon-1)^2(1-p)^2}
{(e^\varepsilon+3)h_{\{x_0\}}(\theta)},
\qquad
\eta_{\mathrm{treatment}}
=
\frac{(e^\varepsilon-1)^2p^2}
{(e^\varepsilon+3)h_{\{x_3\}}(\theta)}.
\]
The first and fourth singleton summands give, respectively,
\[
0\le\eta_{\mathrm{control}}t^2
\le F_\theta(\alpha_{\mathrm{sing}},t),
\qquad
0\le\eta_{\mathrm{treatment}}(t+1)^2
\le F_\theta(\alpha_{\mathrm{sing}},t).
\]
At a minimizing coordinate \(t_{\mathrm{sing}}\), a zero objective
would force both \(t_{\mathrm{sing}}=0\) and
\(t_{\mathrm{sing}}+1=0\). Therefore
\[
0<
F_\theta(\alpha_{\mathrm{sing}},t_{\mathrm{sing}})
=
\mathscr P(\alpha_{\mathrm{sing}})
\le J^*(\theta,p,\varepsilon).
\]

For completeness, the output information matrix is positive
semidefinite. Define its dominating measure and density quantities
by
\[
\nu_Q=\sum_{j=0}^3Q(\,\cdot\mid x_j),
\qquad
d_{Q,j}=\frac{dQ(\,\cdot\mid x_j)}{d\nu_Q},
\]
\[
f_{\theta,Q}(z)=\sum_{j=0}^3\pi_{\theta,j}d_{Q,j}(z),
\qquad
\dot f_{\theta,Q}(z)=
\sum_{j=0}^3(\nabla_\theta\pi_{\theta,j})d_{Q,j}(z).
\]
The information integrand is assigned zero when its denominator
vanishes. Symmetry follows from its outer-product form, and, for
every \(w\in\mathbb R^2\),
\[
w^\top I_\theta(Q)w
=
\int
\frac{(w^\top\dot f_{\theta,Q}(z))^2}
{f_{\theta,Q}(z)}\,d\nu_Q(z)\ge0.
\]
Here \(f_{\theta,Q}\ge0\); the interior parameters and the finite
record alphabet justify the finite-sum and integral operations.

We apply the contrast-variance bound directly, including its range
case split. If \(c\notin\operatorname{range}(I_\theta(Q))\), the
variance is \(+\infty\), which proves the required inequality.
Otherwise choose \(w_Q\in\mathbb R^2\) with
\(I_\theta(Q)w_Q=c\), and define
\[
\ell_Q=c^\top w_Q=w_Q^\top I_\theta(Q)w_Q.
\]
Positive semidefiniteness gives \(\ell_Q\ge0\). Equality would imply
\(I_\theta(Q)w_Q=0\), contradicting \(c\ne0\), so \(\ell_Q>0\).
Any other solution \(\widetilde w_Q\in\mathbb R^2\) satisfies
\[
c^\top\widetilde w_Q
=
w_Q^\top I_\theta(Q)\widetilde w_Q
=
w_Q^\top c
=\ell_Q.
\]
Thus the infimum in the contrast-variance definition equals
\(\ell_Q\).

Define the vector
\[
y_Q=\ell_Qv(t_Q)-w_Q\in\mathbb R^2.
\]
Using \(v(t_Q)^\top c=1\), expansion yields
\[
0\le y_Q^\top I_\theta(Q)y_Q
=
\ell_Q\left(
\ell_Qv(t_Q)^\top I_\theta(Q)v(t_Q)-1
\right).
\]
Since \(\ell_Q>0\), the intermediate bound above implies
\[
1
\le
\ell_Qv(t_Q)^\top I_\theta(Q)v(t_Q)
\le
\ell_QJ^*(\theta,p,\varepsilon).
\]
Taking the reciprocal of the positive optimized information gives
\[
V^*(\theta,p,\varepsilon)
=\frac{1}{J^*(\theta,p,\varepsilon)}
\le\ell_Q
=c^\top I_\theta(Q)^+c.
\]
% lean: CausalSmith.Stat.LdpAteEfficiencySurface.staircase_refinement_interior
% lean: CausalSmith.Stat.LdpAteEfficiencySurface.Jstar_pos_interior
% lean: CausalSmith.Stat.LdpAteEfficiencySurface.channelFisherInfo_posSemidef
% lean: CausalSmith.Stat.LdpAteEfficiencySurface.reciprocal_le_contrastVariance

\item Choose an optimal feasible weight
\(\alpha_{\mathrm{opt}}\in\mathcal A_\varepsilon\) and a minimizing
coordinate \(t_{\mathrm{opt}}\in\mathbb R\). Optimality of \(\alpha_{\mathrm{opt}}\) and minimality of
\(t_{\mathrm{opt}}\) give
\[
F_\theta(\alpha_{\mathrm{opt}},t_{\mathrm{opt}})
=
\mathscr P(\alpha_{\mathrm{opt}})
=
J^*(\theta,p,\varepsilon).
\]
Define
\[
I_{\mathrm{opt}}=I_\theta(\alpha_{\mathrm{opt}}),
\qquad
j_{\mathrm{opt}}
=
v(t_{\mathrm{opt}})^\top I_{\mathrm{opt}}v(t_{\mathrm{opt}})
=
J^*(\theta,p,\varepsilon)>0.
\]
The matrix is symmetric, and
\[
w^\top I_{\mathrm{opt}}w
=
\sum_{S\in\mathcal S}
(\alpha_{\mathrm{opt}})_S
\frac{(g_S^\top w)^2}{h_S(\theta)}
\ge0
\qquad(w\in\mathbb R^2).
\]
Moreover, the minimizing property becomes
\[
v(t_{\mathrm{opt}})^\top I_{\mathrm{opt}}v(t_{\mathrm{opt}})
\le
v(u)^\top I_{\mathrm{opt}}v(u)
\qquad(u\in\mathbb R).
\]

To compute its contrast variance, define
\[
\mathbf 1_2=\begin{pmatrix}1\\1\end{pmatrix},
\qquad
a_{\mathrm{shift}}=\mathbf 1_2^\top I_{\mathrm{opt}}\mathbf 1_2,
\qquad
b_{\mathrm{shift}}
=2\mathbf 1_2^\top I_{\mathrm{opt}}v(t_{\mathrm{opt}}).
\]
Since \(v(t_{\mathrm{opt}}+\zeta)
=v(t_{\mathrm{opt}})+\zeta\mathbf 1_2\), minimality implies
\[
0\le
\zeta b_{\mathrm{shift}}+\zeta^2a_{\mathrm{shift}}
\qquad(\zeta\in\mathbb R).
\]
This forces \(b_{\mathrm{shift}}=0\). Indeed, if it were nonzero,
define
\[
D_{\mathrm{shift}}=|a_{\mathrm{shift}}|+1>0,
\qquad
\zeta_{\mathrm{shift}}
=-\frac{b_{\mathrm{shift}}}{2D_{\mathrm{shift}}}.
\]
Then \(a_{\mathrm{shift}}\le D_{\mathrm{shift}}\) gives
\[
\zeta_{\mathrm{shift}}b_{\mathrm{shift}}
+\zeta_{\mathrm{shift}}^2a_{\mathrm{shift}}
=
-\frac{b_{\mathrm{shift}}^2}{2D_{\mathrm{shift}}}
+\frac{a_{\mathrm{shift}}b_{\mathrm{shift}}^2}
{4D_{\mathrm{shift}}^2}
\le
-\frac{b_{\mathrm{shift}}^2}{4D_{\mathrm{shift}}}<0,
\]
a contradiction.

Thus the two coordinates of
\(I_{\mathrm{opt}}v(t_{\mathrm{opt}})\) sum to zero. Combining this
with the definition of \(j_{\mathrm{opt}}\) gives
\[
I_{\mathrm{opt}}v(t_{\mathrm{opt}})
=
\begin{pmatrix}-j_{\mathrm{opt}}\\j_{\mathrm{opt}}\end{pmatrix}
=j_{\mathrm{opt}}c.
\]
Consequently, the vector
\[
w_{\mathrm{opt}}
=\frac{v(t_{\mathrm{opt}})}{j_{\mathrm{opt}}}
\]
satisfies
\[
I_{\mathrm{opt}}w_{\mathrm{opt}}=c,
\qquad
c^\top w_{\mathrm{opt}}=\frac1{j_{\mathrm{opt}}}.
\]
For every \(w\in\mathbb R^2\) satisfying \(I_{\mathrm{opt}}w=c\),
symmetry gives
\[
c^\top w
=w_{\mathrm{opt}}^\top I_{\mathrm{opt}}w
=w_{\mathrm{opt}}^\top c
=\frac1{j_{\mathrm{opt}}}.
\]
Thus the defining infimum gives
\[
c^\top I_\theta(\alpha_{\mathrm{opt}})^+c
=
\frac1{j_{\mathrm{opt}}}
=
V^*(\theta,p,\varepsilon).
\]

It remains to identify the corresponding channel information.
Feasibility makes each row of \(Q_{\alpha_{\mathrm{opt}}}\) a
probability distribution. For all \(S\in\mathcal S\) and
\(x,x'\in\mathcal X\),
\[
b_S(x)\le e^\varepsilon b_S(x'),
\qquad
(\alpha_{\mathrm{opt}})_Sb_S(x)
\le
e^\varepsilon(\alpha_{\mathrm{opt}})_Sb_S(x').
\]
Summing over any subset of the finite output alphabet gives the
required privacy inequality.

For \(S\in\mathcal S\), define
\[
r_S=\sum_{j=0}^3b_S(x_j)>0.
\]
The dominating measure and conditional densities of this channel
satisfy
\[
\nu_{Q_{\alpha_{\mathrm{opt}}}}(\{S\})
=(\alpha_{\mathrm{opt}})_Sr_S,
\qquad
d_{Q_{\alpha_{\mathrm{opt}}},j}(S)
=
\begin{cases}
b_S(x_j)/r_S,&(\alpha_{\mathrm{opt}})_S>0,\\
0,&(\alpha_{\mathrm{opt}})_S=0.
\end{cases}
\]
The zero-weight outputs have zero dominating mass. At every
positive-weight output, the output density and its derivative are
\[
f_{\theta,Q_{\alpha_{\mathrm{opt}}}}(S)
=\frac{h_S(\theta)}{r_S},
\qquad
\dot f_{\theta,Q_{\alpha_{\mathrm{opt}}}}(S)
=\frac{g_S}{r_S}.
\]
Hence integration over the finite alphabet gives
\[
\begin{aligned}
I_\theta(Q_{\alpha_{\mathrm{opt}}})
&=
\sum_{\substack{S\in\mathcal S\\(\alpha_{\mathrm{opt}})_S>0}}
(\alpha_{\mathrm{opt}})_Sr_S
\frac{g_Sg_S^\top}{h_S(\theta)r_S}\\
&=
\sum_{S\in\mathcal S}
(\alpha_{\mathrm{opt}})_S
\frac{g_Sg_S^\top}{h_S(\theta)}
=
I_\theta(\alpha_{\mathrm{opt}}).
\end{aligned}
\]
Together with the preceding variance calculation, this proves
\[
c^\top I_\theta(Q_{\alpha_{\mathrm{opt}}})^+c
=
c^\top I_\theta(\alpha_{\mathrm{opt}})^+c
=
V^*(\theta,p,\varepsilon).
\]
The universal lower bound and this admissible attaining channel
then yield
\[
\inf_Qc^\top I_\theta(Q)^+c=V^*(\theta,p,\varepsilon),
\]
with the infimum taken over the channel class in the statement.
% lean: CausalSmith.Stat.LdpAteEfficiencySurface.Jstar_exists_optimal_weight
% lean: CausalSmith.Stat.LdpAteEfficiencySurface.contrastVariance_eq_reciprocal_of_minimizer
% lean: CausalSmith.Stat.LdpAteEfficiencySurface.staircaseChannel_stationaryLDP
% lean: CausalSmith.Stat.LdpAteEfficiencySurface.channelFisherInfo_staircase_eq_informationMatrix

\item Finally, set
\[
t^*=t_{\mathrm{env}}\in\mathbb R.
\]
Since \(t_{\mathrm{env}}\) minimizes \(\mathscr E\), the minimax identity
gives
\[
\max_{\alpha\in\mathcal A_\varepsilon}
F_\theta(\alpha,t^*)
=
\mathscr E(t_{\mathrm{env}})
=
\inf_{t\in\mathbb R}\mathscr E(t)
=
J^*(\theta,p,\varepsilon),
\]
which proves the asserted attainment of the upper-envelope value.
% lean: CausalSmith.Stat.LdpAteEfficiencySurface.upperEnvelope_exists_minimizer
% lean: CausalSmith.Stat.LdpAteEfficiencySurface.finite_oracle
\end{enumerate}
\end{proof}
